\documentclass[11pt]{article}
\usepackage[utf8]{inputenc}
\usepackage[T1]{fontenc}
\usepackage[italian]{babel}
\usepackage{geometry}
\geometry{a4paper, margin=1in}
\usepackage{xcolor}
\usepackage{listings}
\usepackage{hyperref}

\usepackage{tcolorbox}
\tcbuselibrary{breakable}
\begin{document}

\begin{tcolorbox}[colback=blue!5!white,colframe=blue!75!black,title=\textbf{DOMANDA}, breakable]
Esercizio
\begin{itemize}
    \item Modificare il file \texttt{serverHTTP.c} in modo che, invece di restituire sempre la solita pagina web di prova, restituisca una delle pagine html usate dal lucido 9 in poi scelta attraverso l'uso del browser.
    \item Suggerimenti:
    \begin{itemize}
        \item Provare a fare la nuova richiesta col browser usando il \texttt{serverHTTP.c} in modo da vedere la richiesta HTTP per capire dove si trova la stringa con il nome del file nella richiesta che fa il browser (aiutarsi anche con Wireshark)
    \end{itemize}
    \item Cosa devo scrivere nella barra del browser?
    \begin{itemize}
        \item Capire come recuperare la stringa con il nome del file dalla richiesta HTTP (si veda \texttt{esempio-parser.c})
        \item Per costruire la risposta riciclare parte del codice usato nell'esercizio del trasferimento di file
    \end{itemize}
    \item Prova finale: cosa succede se chiedo al server di restituire i file \texttt{form-get.html} e \texttt{form-post.html}?
\end{itemize}
\end{tcolorbox}

\begin{tcolorbox}[colback=green!5!white,colframe=green!75!black,title=\textbf{RISPOSTA}, breakable]
\textbf{Soluzione Esercizio 25:} Web server e file dinamici (URL)

\textbf{Operazioni Preliminari e Modifica del Codice}\\
Per questo esercizio ho creato la cartella e ho modificato in modo sostanziale il file "\texttt{serverHTTP.c}":
\begin{enumerate}
    \item Utilizzando la funzione \texttt{strtok} (come suggerito nel file \texttt{esempio-parser.c}), ho estratto la seconda parola della prima riga della richiesta \texttt{HTTP}, ovvero l'URI (es: "/prova.html").
    \item Ho manipolato questa stringa per rimuovere lo slash "/" iniziale e ho inserito una logica che ripulisce la stringa da eventuali "Query String" (es. rimuovendo tutto ciò che c'è da un eventuale punto interrogativo "?" in poi).
    \item Utilizzando la funzione \texttt{fopen}, il server tenta di aprire il file. Se il file esiste, invia prima di tutto gli Header di successo (\texttt{HTTP/1.1 200 OK}) e poi, tramite un ciclo \texttt{fgetc}, legge e trasmette al client l'intero contenuto del file. Se il file non esiste, il server restituisce una pagina d'errore dinamica (\texttt{HTTP/1.1 404 Not Found}).
\end{enumerate}

\textbf{Risposta alla Domanda 1: "Cosa devo scrivere nella barra del browser?"}\\
Per poter richiedere una pagina specifica presente nella cartella Esempi-web (che abbiamo copiato all'interno dell'esercizio), è necessario digitare l'indirizzo base del server aggiungendo lo slash e il nome esatto del file.\\
\textbf{Ad esempio:} \texttt{http://localhost:8000/prova.html} oppure \texttt{http://localhost:8000/javascript.html}.

\textbf{Risposta alla Domanda 2: "Cosa succede se chiedo al server di restituire i file form-get.html e form-post.html?"}\\
Scrivendo l'indirizzo \texttt{http://localhost:8000/form-get.html}:
\begin{enumerate}
    \item Il browser effettua una richiesta \texttt{HTTP} GET al nostro server C.
    \item Il server C individua il file sul disco, ne legge il contenuto testuale (che è formato da tag \texttt{<form>}, \texttt{<input>}, ecc.) e lo rispedisce intatto al browser.
    \item Il browser, ricevendo il file, lo \textbf{renderizza graficamente}, mostrando all'utente la classica maschera per inserire nome e cognome.
    \item Se l'utente compila la form e preme il pulsante "Submit" (e l'azione della form è \texttt{action="/action"} come in originale), il browser avvierà autonomamente una \textbf{seconda} richiesta \texttt{HTTP} (\texttt{GET /action?fname=...}) verso il nostro server C. Poiché il file "action" non esiste, in questo caso il nostro server risponderà con un Errore 404 Not Found!
\end{enumerate}

\textbf{Questo ci permette di capire a fondo come comunicano Front-end e Back-end:} il server prima serve l'interfaccia utente (HTML) su richiesta del client, e successivamente il client usa quell'interfaccia per inviare comandi (dati e parametri) allo stesso server.

\textbf{Comandi Utili per testare l'Esercizio:}
\begin{lstlisting}[language=bash, basicstyle=\ttfamily\small]
# Compila ed esegui il server
gcc serverHTTP.c network.c -o serverHTTP
./serverHTTP

# Richiedi la pagina prova.html da riga di comando per testarlo (oppure apri l'URL sul browser):
curl -v http://localhost:8000/prova.html
curl -v http://localhost:8000/file_inesistente.txt
\end{lstlisting}
\end{tcolorbox}

\end{document}
